\documentclass[10pt]{beamer}
% Adapted from the supervisor-supplied BCIT Beamer template.
% Original 4:3 aspect ratio, theme, colours and logo placement retained.
\usetheme[progressbar=head,numbering=fraction,sectionpage=none]{metropolis}
\usepackage{appendixnumberbeamer}
\usepackage{booktabs,array,amsmath,graphicx,tikz}
\newcolumntype{P}[1]{>{\raggedright\arraybackslash}p{#1}}
\usetikzlibrary{positioning,arrows.meta}
\definecolor{bcitblue}{HTML}{003C71}
\definecolor{bcityellow}{HTML}{FFEA2E}
\definecolor{bcitlight}{HTML}{F8F9FF}
\definecolor{bcitgray}{HTML}{6E6E6E}
\setbeamercolor{palette primary}{bg=bcitblue,fg=bcitlight}
\setbeamercolor{progress bar}{fg=bcityellow,bg=bcitblue!30}
\setbeamercolor{progress bar in head/foot}{fg=bcityellow,bg=bcitblue!30}
\setbeamercolor{frametitle}{bg=bcitblue,fg=bcitlight}
\setbeamercolor{block title}{bg=bcitblue!90,fg=bcitlight}
\setbeamercolor{block body}{bg=bcitblue!10,fg=black}
\setbeamercolor{block title alerted}{bg=bcityellow!90,fg=bcitblue}
\setbeamercolor{block body alerted}{bg=bcityellow!10,fg=black}
\setbeamercolor{item}{fg=bcitblue}
\setbeamercolor{itemize item}{fg=bcitblue}
\setbeamertemplate{blocks}[rounded]
% Fit the template title content inside the plain frame's available height.
\setbeamertemplate{title page}{
  \begin{minipage}[b][\textheight]{\textwidth}
  \vfill
  \usebeamertemplate*{title}
  \usebeamertemplate*{subtitle}
  \usebeamertemplate*{title separator}
  \usebeamertemplate*{author}
  \usebeamertemplate*{date}
  \usebeamertemplate*{institute}
  \vfill
  \end{minipage}
}
\setbeamertemplate{frame footer}{\raisebox{-0.1cm}{\includegraphics[height=0.4cm]{bcit_logo.png}}\hspace{0.4cm}\tiny\insertshorttitle}
\makeatletter
\setlength{\metropolis@progressonsectionpage@linewidth}{2pt}
\setlength{\metropolis@progressinheadfoot@linewidth}{2pt}
\makeatother
\newcommand{\src}[1]{\vfill{\tiny\color{bcitgray}#1}}
\newcommand{\cat}[1]{\texttt{category\_#1}}
\title[COMP 9500: Natural-language controls]{Natural-Language Controls for\\Short-Video Recommenders}
\subtitle{An offline study of controllability}
\author{Vibhor Malik}
\date{Supervisor review draft: 4 October 2026, Vancouver time}
\institute{COMP 9500, M.Sc. Research Project\\British Columbia Institute of Technology\\Supervisor: Borna Noureddin}
\begin{document}

{\setbeamertemplate{frame footer}{}
\setbeamertemplate{background}{\begin{tikzpicture}[remember picture,overlay]
\node[anchor=north east,xshift=-0.5cm,yshift=-0.5cm] at (current page.north east){\includegraphics[height=0.8cm]{bcit_logo.png}};
\end{tikzpicture}}
\begin{frame}[plain]

\maketitle
\end{frame}}

\section{Research question}
\begin{frame}{Current research position}
\begin{itemize}
\item Experiment 1 establishes category enforcement under a fixed policy.
\item The LLM adds no ranking advantage on the templated task, and boosting costs predictive quality.
\item Experiment 2 tests richer intents, but its reserved evaluation remains blocked.
\end{itemize}
\begin{alertblock}{Claim boundary}
No demonstrated expressiveness advantage, accuracy gain or improvement in human experience.
\end{alertblock}
\src{Evidence: Experiment 1 summaries and Experiment 2 status record, archived 29 September 2026.}
\end{frame}

\begin{frame}{Research question}
Can one natural-language interaction satisfy a richer intended change than one item reaction?
\vspace{0.35cm}
\begin{itemize}
\item A compound request can specify several categories, strength and exceptions.
\item The LLM must also outperform a rule parser to establish its own contribution.
\item Offline compliance and historical prediction measure different outcomes.
\end{itemize}
\vspace{0.2cm}
The claim that language \emph{feels more natural} requires human evidence beyond this study.
\end{frame}

\section{Related work}
\begin{frame}{Prior work supports a control question}
\small
\begin{tabular}{@{}P{2.1cm}P{8.0cm}@{}}
\toprule
Source & Relevant evidence\\\midrule
Mozilla, 2022 & Dislike prevented 12\% of unwanted recommendations and Not interested 11\%, relative to no feedback.\\[0.15cm]
Ramos et al., 2024 & Positive profile edits raised target Coverage@10 without further fine-tuning. Amazon-MT RMSE: 0.941 versus MF 0.925, where lower is better.\\[0.15cm]
Sanner et al., 2023 & Language preferences were competitive in a near cold-start movie setting.\\[0.15cm]
Chen and Pu, 2012 & Critiquing systems already studied richer explicit preference feedback.\\\bottomrule
\end{tabular}\par
\src{Historical results in different settings. Full references and source links: backup slide and speaker notes.}
\end{frame}

\section{Experiment 1}
\begin{frame}{Translation and enforcement are separate}
\begin{tabular}{@{}P{2.2cm}P{7.9cm}@{}}
\toprule
Stage & Role\\\midrule
Recommender & Frozen SASRec-style scores for the candidate catalogue.\\[0.18cm]
Translator & Llama 3.1 8B or rules compile language into a validated command.\\[0.18cm]
Enforcer & Shared deterministic reranking applies the command.\\[0.18cm]
Explanation & Verified facts describe the displayed item and control.\\\bottomrule
\end{tabular}\par
\vspace{0.35cm}
A correct hard exclusion comes from the enforcement code. It does not demonstrate language understanding.
\end{frame}

\begin{frame}{Data and information boundaries}
\begin{itemize}
\item KuaiRec cohort: \textbf{865 users}, \textbf{2,593 candidates}, 31 opaque category identifiers.
\item Positive event: watch ratio $>2$, an engagement proxy.
\item Global chronological windows: nominal 70\% training, 15\% request history, 15\% held-out evaluation.
\item Requests use training and request history only.
\item Quality target: first eligible held-out positive item.
\end{itemize}
\vspace{0.2cm}
Synthetic requests and numeric category labels limit claims about real preferences and language.
\src{Filtered cohort counts, not the upstream KuaiRec matrix size. Frozen candidate policy permits seen items.}
\end{frame}

\begin{frame}{Five matched configurations}
\small
\begin{tabular}{@{}cP{8.9cm}@{}}
\toprule
ID & Mechanism\\\midrule
A & One simulated reaction changes one historical exemplar item.\\[0.15cm]
B & A plus display-only explanations.\\[0.15cm]
C & Free text compiles into a category control.\\[0.15cm]
D & An editable preference sentence compiles into a category control.\\[0.15cm]
E & Text, profile and explanations, with duplicate commands removed.\\\bottomrule
\end{tabular}\par
\vspace{0.25cm}
One boost and one mute request per user, under seeds 42, 43 and 44. The item-feedback proxy is an explicit design choice.
\src{A is not a reconstruction of a commercial feedback algorithm. B must rank identically to A.}
\end{frame}

\begin{frame}{Parser results on synthetic language}
\begin{tabular}{@{}lrr@{}}
\toprule
Language material & Llama correct & Meaning\\\midrule
Templated inputs & 120/120 & Canonical task\\
Broader synthetic set & 79/100 & Below 85\% target\\
Adversarial subset & 2/10 & Major weakness\\\bottomrule
\end{tabular}\par
\vspace{0.3cm}
\begin{itemize}
\item 79\% exceeds the 70\% floor. It misses the 85\% target.
\item The rule parser scores 81/100 on the same broader set.
\item The 120 canonical wording patterns informed development.
\end{itemize}
\src{The adversarial subset belongs to the 100 cases. These finite synthetic sets do not establish general language reliability.}
\end{frame}

\begin{frame}{Category enforcement works under the policy}
\begin{tabular}{@{}lrr@{}}
\toprule
Top-10 exposure & A & C, D and E\\\midrule
Boost target share, TCP & 11.6\% & 79.2\%\\[0.15cm]
Muted target share, TCER & 6.1\% & 0.0\%\\\bottomrule
\end{tabular}\par
\vspace{0.4cm}
\begin{itemize}
\item Hard exclusion guarantees zero exposure when parsed correctly.
\item Category-wide boosts affect more items than A's single-item adjustment.
\item Rule and LLM paths produce identical feeds for this task.
\end{itemize}
\src{Averages over 865 users and three seeds. This is enforcement evidence, not an LLM advantage.}
\end{frame}

\begin{frame}{Steering has a predictive quality cost}
\begin{tabular}{@{}lrr@{}}
\toprule
Boost requests & A & C\\\midrule
NDCG@10 & 0.03012 & 0.01637\\\bottomrule
\end{tabular}\par
\vspace{0.35cm}
\textbf{45.7\% relative decrease.}
\vspace{0.3cm}

Paired mean change: $-0.01375$\par
95\% user-bootstrap interval: $[-0.02086,-0.00655]$.
\vspace{0.3cm}

This post-hoc secondary contrast measures agreement with a historical next item. It does not measure user dissatisfaction after steering.
\src{Post-hoc exploratory interval, outside the frozen primary BH family.}
\end{frame}

\begin{frame}{A weak backbone limits the quality claim}
\begin{tabular}{@{}lr@{}}
\toprule
Uncontrolled ranking & HR@1\\\midrule
SASRec, seed 42 & 0.000\%\\
SASRec, seed 43 & 0.347\%\\
SASRec, seed 44 & 0.462\%\\
Recency popularity & 7.977\%\\\bottomrule
\end{tabular}\par
\vspace{0.3cm}
Recency pools all users' request history. The model's weights use the older training window. This comparison is descriptive.
\vspace{0.2cm}
The absolute 10 percentage point non-inferiority margin is uninformative at this baseline.
\src{Post-hoc uncontrolled baselines. A pooled HR@1 interval includes zero. No accuracy gain is established.}
\end{frame}

\begin{frame}{Experiment 1 conclusion}
\begin{enumerate}
\item Deterministic category control works: TCP 11.6\% to 79.2\%, TCER 6.1\% to zero.
\item Llama parsed 120/120 templated inputs, 79/100 broader synthetic inputs and 2/10 adversarial cases.
\item All 10/10 explanation fact selections equal the template.
\item Boost NDCG@10 falls 45.7\%. Non-inferiority is uninformative.
\item Backbone HR@1 is 0.0--0.46\%, below recency popularity's 8.0\%.
\end{enumerate}
\src{The study establishes bounded enforcement and its limitations. It does not establish an LLM advantage or useful quality preservation.}
\end{frame}

\section{Experiment 2}
\begin{frame}{Experiment 2: richer intent per interaction}
\begin{description}
\item[H1] Higher induced intent satisfaction than both the one-item proxy and the rule parser.
\item[H2] More intended-direction rank movement than one item reaction, with one interaction each.
\item[H3] Paired NDCG@10 and HR@1 changes against A and recency popularity, with intervals.
\end{description}
\vspace{0.25cm}
H3 is descriptive. It has no pass threshold and no non-inferiority claim.
\src{Prespecified protocol dated 29 September 2026. H1 and H2 remain unassessed.}
\end{frame}

\begin{frame}{Four classes of intent}
\small
\begin{tabular}{@{}P{2.0cm}P{4.0cm}P{3.7cm}@{}}
\toprule
Class & Illustrative request & Top-10 requirement\\\midrule
Compound & More X and less Y & $n(X)\geq2$, $n(Y)\leq1$\\[0.18cm]
Graded & A bit more X & $2\leq n(X)\leq4$\\
 & Much less X / none of X & $n(X)\leq1$ / $n(X)=0$\\[0.18cm]
Floor & Less X, but not zero & $1\leq n(X)\leq2$\\[0.18cm]
Conditional & More X only if also Y & $n(X\cap Y)\geq2$ and no increase in $X\setminus Y$ count\\\bottomrule
\end{tabular}\par
\vspace{0.25cm}
Success also requires a full list and positive intended movement. Bands are operational conventions.
\src{X and Y denote opaque categories. Conditional means the same item carries both labels. These are illustrative wording forms, not reserved cases.}
\end{frame}

\begin{frame}{A matched interaction budget}
\begin{itemize}
\item A gets one positive or negative reaction on one historical exemplar.
\item Rule and LLM get the complete utterance once and share the same ranker.
\item Assignments use only earlier history and a committed SHA-256 scheme.
\item Channels share permitted pre-test information. Recency and SASRec use that information differently.
\end{itemize}
\vspace{0.3cm}
The budget matches interaction count. It does not equalize semantic content, user effort or every possible feedback algorithm.
\end{frame}

\begin{frame}{Reserved evaluation procedure}
\begin{enumerate}
\item Preserve the 48 development and 100 reserved cases.
\item Retain the parser commitment and record prior exposure.
\item Obtain the student's independent annotation, agreement and blinded adjudication.
\item Complete the freeze and verify the Llama runtime.
\item Open reserved evaluation once. Report parser accuracy before ranking results.
\end{enumerate}
\vspace{0.2cm}
AI annotation cannot supply the required independent student pass.
\src{Separate development, annotation and evaluation accesses have distinct receipts. No tuning on reserved outcomes.}
\end{frame}

\begin{frame}{Development evidence is not a test result}
\begin{tabular}{@{}lrr@{}}
\toprule
Development version & Rule parser & Llama\\\midrule
v1 & 48/48 & 24/48\\
v2, after development & 48/48 & 48/48\\\bottomrule
\end{tabular}\par
\vspace{0.3cm}
\begin{itemize}
\item Prompt and wording changed during development.
\item v1 includes 10 operational errors for Llama.
\item The 100 reserved parser outcomes remain unmeasured.
\item Recorded reserved ranking rows: \textbf{0}.
\end{itemize}
\src{Archived development runs dated 29 September 2026. Equal development scores establish no LLM advantage.}
\end{frame}

\begin{frame}{Prespecified analysis}
\begin{itemize}
\item Unit: user. Average three seeds within each user.
\item H1: four classes and pooled, against A and rules.
\item H2: pooled intended movement against A.
\item One family of \textbf{11} Wilcoxon comparisons, BH correction.
\item User-bootstrap intervals from \textbf{5,000} resamples.
\item Retain impossible and already-satisfied cases with explicit flags.
\end{itemize}
\vspace{0.2cm}
H1 needs both positive pooled contrasts with $q<0.05$. H2 needs its positive pooled contrast with $q<0.05$.
\src{Signed-rank symmetry, tied/sparse differences and BH dependence assumptions remain limitations.}
\end{frame}

\section{Completion and discussion}
\begin{frame}{Current completion limits}
\begin{tabular}{@{}P{3.3cm}P{6.7cm}@{}}
\toprule
Component & Status\\\midrule
Experiment 1 & Finished and frozen.\\[0.13cm]
Experiment 2 H1 / H2 & Not assessed. No pass or failure claim.\\[0.13cm]
Experiment 2 H3 & Not run. No quality conclusion.\\[0.13cm]
Student annotation & Required genuine human pass is missing.\\[0.13cm]
Local Llama runtime & Unavailable at the 3 October check.\\[0.13cm]
Windows and Colab & Full research execution unverified.\\\bottomrule
\end{tabular}\par
\src{Software verification does not complete the blocked experiment. No rule output substitutes for Llama.}
\end{frame}

\begin{frame}{Contribution and publication direction}
\textbf{Supported contribution}
\begin{itemize}
\item An inspectable implementation separating translation from enforcement.
\item No observed ranking benefit from LLM parsing on the tested templates, with an explicit quality cost.
\item A prospective protocol for testing richer feedback against both relevant comparators.
\end{itemize}
\vspace{0.3cm}
\textbf{Work needed for a stronger paper}
\begin{itemize}
\item Complete the reserved study and report its actual outcome.
\item Justify the intent measures and resolve scope with the supervisor.
\item Evaluate broader language and human experience in a separately approved extension.
\end{itemize}
\end{frame}

\begin{frame}{Proposed 135-hour completion plan}
\small
\begin{tabular}{@{}P{2.2cm}P{7.7cm}@{}}
\toprule
Proposed period & Main checkpoint\\\midrule
5--23 Oct & Scope, sources, annotation and runtime. Readiness decision by 23 October.\\[0.18cm]
26 Oct--8 Nov & Reserved evaluation 26 Oct--1 Nov if ready. Analysis 2--8 November.\\[0.18cm]
9--27 Nov & Reproduction 9--15 Nov. Cold review 16--22 Nov. Complete paper draft by 27 Nov.\\[0.18cm]
30 Nov--18 Dec & Supervisor revisions, repository handoff, rehearsal and final review.\\\bottomrule
\end{tabular}\par
\vspace{0.25cm}
11 remaining calendar weeks, \textbf{135 prospective hours}. Exact submission date requires confirmation.
\src{Weekly check-in: Wednesday 10:00 am. Update by Tuesday 10:00 am. Meeting notes follow each meeting. No completed hours are claimed here.}
\end{frame}

\begin{frame}{Supervisor discussion}
\begin{enumerate}
\item Confirm the revised scope, approval record and report format.
\item Resolve the original non-inferiority margin wording.
\item Agree how prior exposure affects the independent annotation pass.
\item Confirm the final date and repository access route.
\end{enumerate}
\vspace{0.5cm}
The immediate conclusion is bounded: Experiment 1 establishes enforcement, while Experiment 2's expressiveness claim remains untested.\par
\vspace{0.3cm}
\textbf{Questions and discussion: 10--15 minutes.}
\end{frame}

\pdfstringdefDisableCommands{\def\translate#1{#1}}
\appendix
\begin{frame}{Backup: metric definitions}
\begin{description}
\item[TCP@10] Target-category items divided by ten for boost requests.
\item[TCER@10] Target-category items divided by ten for mute requests.
\item[HR@1] One if the first item equals the held-out target, otherwise zero.
\item[NDCG@10] For one relevant item at rank $r\leq10$, $1/\log_2(r+1)$. Zero if absent.
\end{description}
\vspace{0.3cm}
Both exposure metrics have the same arithmetic and opposite desired directions. Historical quality does not directly measure preference satisfaction after intervention.
\end{frame}

\begin{frame}{Backup: backbone and control settings}
\begin{tabular}{@{}P{3.8cm}P{6.0cm}@{}}
\toprule
Backbone & Compact SASRec-style implementation\\\midrule
Hidden size / context & 32 / 50 items\\
Attention & Two blocks, one head\\
Training & 20 epochs, Adam, learning rate 0.001\\
Batch / dropout & 64 / 0.2\\
Training seeds & 42, 43, 44\\
Experiment 1 control & Normalized score +0.25, or hard mute\\
Experiment 2 strengths & 0.10, 0.25, 0.50\\
Language backend & Llama 3.1 8B, temperature 0\\\bottomrule
\end{tabular}\par
\src{The backbone is independently implemented. The study does not claim an exact SASRec paper reproduction.}
\end{frame}

\begin{frame}{Backup: why non-inferiority is uninformative}
\small
\begin{tabular}{@{}P{6.6cm}r@{}}
\toprule
Pooled Experiment 1 quantity & Value\\\midrule
A HR@1 baseline & 0.2505\%\\
C minus A, mean & +0.193 pp\\
Paired 95\% interval & $[-0.058,+0.501]$ pp\\
Simultaneous lower bound & $-0.07707$ pp\\
Absolute proposed threshold & $-10$ pp\\
Relative 10\% loss threshold & $-0.02505$ pp\\\bottomrule
\end{tabular}\par
\vspace{0.3cm}
The absolute reading numerically passes while allowing the entire small baseline to disappear. The relative reading fails. Only eight of 865 users have nonzero paired HR@1 differences.
\src{Both operations pooled, distinct from the boost-only NDCG contrast. pp means percentage points.}
\end{frame}

\begin{frame}{Backup: intended-direction movement}
\small
\[
D(S)=\frac{\sum_{r=1}^{10}\mathbf{1}[i_r\in S]/\log_2(r+1)}{\sum_{r=1}^{10}1/\log_2(r+1)}
\]
$\Delta D$ is the channel value minus its uncontrolled value.
\begin{tabular}{@{}P{2.7cm}P{7.0cm}@{}}
\toprule
Class & Signed movement $M$\\\midrule
Compound & $\min(\Delta D(X),-\Delta D(Y))$\\
Boost / reduction & $\Delta D(X)$ / $-\Delta D(X)$\\
Floor & $-\Delta D(X)$, zero if the floor is violated\\
Conditional & $\Delta D(X\cap Y)-\max(0,\Delta D(X\setminus Y))$\\\bottomrule
\end{tabular}\par
\vspace{0.25cm}
Induced satisfaction requires a full list, the class band and $M>10^{-12}$. This operational measure is not Shannon information.
\end{frame}

\begin{frame}{Backup: evidence and responsible authorship}
\small
\begin{itemize}
\item E1 metrics: \texttt{results/llm/seed\_*/request\_metrics.csv}.
\item E1 parser: \texttt{results/parser\_ollama\_test/}.
\item Baselines and paired analysis: \texttt{results/audit\_revision/}.
\item E2 commitments and status: \texttt{results/experiment2/}.
\item E2 protocol: \texttt{docs/experiment2-protocol.md}.
\end{itemize}
\vspace{0.3cm}
AI assisted implementation and documentation. Personal reading, independent annotation and manual execution should be claimed only when completed and recorded by the student.
\src{The current demo illustrates Experiment 1. Inspecting source files is not an execution record.}
\end{frame}

\begin{frame}{Backup: selected references}
\footnotesize
\begin{itemize}
\item Ricks and McCrosky (Mozilla), 2022. \emph{Does this button work? Investigating YouTube's ineffective user controls.}
\item Ramos et al., 2024. \emph{Transparent and Scrutable Recommendations Using Natural Language User Profiles}. arXiv:2402.05810.
\item Sanner et al., RecSys 2023. \emph{Large Language Models are Competitive Near Cold-start Recommenders for Language- and Item-based Preferences}. arXiv:2307.14225.
\item Chen and Pu, 2012. \emph{Critiquing-based recommenders: survey and emerging trends}. UMUAI 22, 125--150.
\item Kang and McAuley, ICDM 2018. \emph{Self-Attentive Sequential Recommendation}.
\item Gao et al., CIKM 2022. \emph{KuaiRec: A Fully-observed Dataset and Insights for Evaluating Recommender Systems}.
\end{itemize}
\src{Complete bibliographic records: paper/references.bib. Direct source links: speaker-notes.md.}
\end{frame}
\end{document}
